\section{Reflection}\label{sec:reflection}

This section addresses the three reflective requirements of the Research Project:
the novelty of the results, their efficacy, and what I learned. I write in the
first person here because the judgements are mine.

\subsection{Novelty}

Three contributions appear to be new, in descending order of confidence.

\textbf{An energy measurement of an agentic AutoML loop for which I found no
direct precedent.} Published
inference-energy work measures single prompts \cite{luccioni2024power};
MLAgentBench \cite{huang2024mlagentbench} measures agent success rates and API
cost but not energy. In a literature search last updated in August~2026, I did
not identify a prior study that measures the energy of a directly comparable
closed propose--evaluate--keep loop. Under this experiment's fixed iteration
budgets, proposal quality determines retained yield rather than how many training
runs are scheduled. The specific finding is a large sparse-proposer energy
advantage with a higher, but statistically uncertain, mean accuracy; it is one
workload and one model pair.

\textbf{Per-device energy attribution as a design choice rather than a model.}
Placing proposer inference and inner-loop training on physically separate GPUs
means GPU~1 and GPU~0 board energy are read from different devices rather than
separated by a calibration model. Whole-session device totals still include idle
and residency energy. Phase-aligned power checks confirm that GPU~0 draws more
during training and GPU~1 during proposing; the looser reconstruction flag itself
is only a coverage check and should not be mistaken for a 1\,\% accuracy test.

\textbf{A methodological result that generalises beyond these subjects.} A keep
threshold must be derived from a measured noise floor; the value I specified
a~priori sat four times below it, at which point ``retained mutation'' is not a
meaningful event. I consider this more portable than the architectural result,
because it applies to any study of this loop shape.

I had intended to claim a second one, that deterministic sampling makes such a
loop degenerate. The follow-up showed that $T=0$ is not sufficient to cause this
under the revised prompt. What
survives is narrower: temperature governs how often the proposer repeats itself.
The available outcome samples do not establish a null effect on retention or
accuracy.

I am less confident about novelty in one respect. The observation that contract
compliance and proposal quality are independent axes is, I suspect, known
informally to practitioners; my contribution is a measurement of it
(0.0\,\% vs.\ 9.6\,\% guard-rejection rates, with the more-rejected arm producing
better recipes) rather than the observation itself.

\subsection{Efficacy: what worked and what did not}

\textbf{What worked.} Calibrating before measuring was decisive. The training
budget, keep threshold and cool-down were each set from measurement, and each
differed materially from the value I had assumed, the budget by a factor of
five, the threshold by a factor of twelve. Had I run the factorial on the assumed
values, the keep/revert decision would have been noise and the study would have
produced a clean-looking null.

The gate structure also worked, in an unglamorous way: gate G1 rejected the
originally proposed model pair on arithmetic before 44\,GB had been downloaded a
second time, and the alignment gate would have caught crossed device indices had
there been any.

\textbf{What did not.} I diagnosed a broken experiment from its first three
sessions and acted on it. Those sessions retained no mutations, I concluded that
sampling at temperature 0 had made the search degenerate, and I changed
temperature and restarted. The reasoning was plausible and the supporting
evidence was real: the proposals were 0.982 similar on average with one pair
identical. It was still overconfident. Later sweeps on both arms showed that
temperature 0 can retain mutations under the revised prompt, and that three
sessions retaining nothing is plausible for the dense arm, which retains nothing
in 42\,\% of its Phase~1 sessions. The
run order is randomised precisely so that the first few runs are not
representative, and I read them as though they were.

I compounded it by changing two things at once. The restart altered sampling and
added the budget facts to the prompt, so even after the sweep I can only say that
temperature 0 is not sufficient under the revised prompt, not what either change contributed on the
dense arm. Eighteen sessions were spent recovering an answer that a
one-change-at-a-time intervention would have given for free.

I also under-powered accuracy relative to energy. An unplanned, functionally
equivalent control (the proposer token caps differed but were non-binding),
executed in two experiments days apart, reproduced GPU-board session energy to
0.20\,\% and test accuracy to only 6.42\,pp. Given that
variance structure, spending the same machine time on more repetitions of fewer
cells would have produced tighter accuracy estimates. I would make that trade
differently if I repeated this.

\subsection{What I learned}

\textbf{An agentic harness is a measuring instrument, and its prompt is part of
the instrument.} Three times the experiment was floored not by any property of
the models but by information the agent could not observe: that the data was
already normalised, that the installed library versions had removed the APIs it
was using, and how few epochs the wall-clock budget actually buys. Each fix was
comment-only and each moved the crash rate substantially. The line I settled on,
supply facts the agent cannot observe and withhold the judgement being measured,
is one I would now apply from the start rather than derive under pressure.

\textbf{Do not diagnose an experiment from its opening runs.} This is the
mistake I would most want to avoid repeating. Randomised run order guarantees
that early sessions are an arbitrary subset, and three of eight cells is not a
sample from which to infer a mechanism. What made it seductive was that I had a
plausible mechanism ready and evidence that fit it; what I lacked was any check
that the evidence discriminated between my explanation and a simpler one.

\textbf{Plausible output is more dangerous than failure.} The failures that cost
me the most were the ones that produced valid-looking data: a profiler faithfully
recording an idle machine because the program it was supposed to wrap had died
silently; an analysis pipeline grouping results by the wrong factor and printing
a plausible table; a run table that a hard stop had left malformed in a way that
only surfaced days later. None raised an error. I have come to trust a crash considerably more than
a clean run, and to check the content of results rather than their shape.

\textbf{Measurement infrastructure is most of the work.} The scientific question
took twenty-four sessions and about ten GPU-hours. Getting to the point where
those sessions meant anything took substantially longer, and most of the
deviations recorded in this report are infrastructural rather than conceptual. I
had not appreciated how much of empirical computer science is instrument-building
before doing it.

\textbf{Reading someone else's analysis script is cheap and worth doing.} Late in
the project I revisited the analysis template from the Green Lab course. Two of
its habits were missing from my pipeline and neither was expensive: plotting the
dependent variable against duration, and against run order. The first turned out
to matter, it separates an efficiency claim from a speed claim, and my
headline result needed that separation to be defensible. The second confirmed
that randomisation had done its job. I had built more sophisticated analysis than
the template and still omitted two of its cheapest and most useful checks, which
suggests I should read comparable work for its routine practice and not only for
its findings.

\textbf{Effect sizes over significance, understood rather than recited.} The
design pre-registered effect sizes as the primary output, which I initially
treated as a convention of the field. The replication above made it concrete:
with agent-side variance of that magnitude at $n=3$, a significance test on
accuracy answers a question about sample size rather than about proposers.
